\documentclass[10pt,a4paper]{amsart}
\usepackage{amsmath,amssymb,amsthm,graphicx}
% Retained minimal preamble for the Stacks Project chapter sites-cohomology.tex (Cohomology on Sites).
% Verbatim excerpt of the chapter preamble: Xy-pic setup (native preamble lines 14-18) and the native
% theorem declarations with the mathematical macros (native preamble lines 154-217). Copied mechanically;
% nothing added, removed or redefined. The wrapper class and the amsmath/amssymb/amsthm/graphicx base
% packages are supplied by the assembler, so the class/externaldocument part of the native preamble is omitted.
% For commutative diagrams we use Xy-pic
\usepackage[all]{xy}

% We use 2cell for 2-commutative diagrams.
\xyoption{2cell}
% Theorem environments.
%
\theoremstyle{plain}
\newtheorem{theorem}[subsection]{Theorem}
\newtheorem{proposition}[subsection]{Proposition}
\newtheorem{lemma}[subsection]{Lemma}

\theoremstyle{definition}
\newtheorem{definition}[subsection]{Definition}
\newtheorem{example}[subsection]{Example}
\newtheorem{exercise}[subsection]{Exercise}
\newtheorem{situation}[subsection]{Situation}

\theoremstyle{remark}
\newtheorem{remark}[subsection]{Remark}
\newtheorem{remarks}[subsection]{Remarks}

\numberwithin{equation}{subsection}

% Macros
%
\def\lim{\mathop{\mathrm{lim}}\nolimits}
\def\colim{\mathop{\mathrm{colim}}\nolimits}
\def\Spec{\mathop{\mathrm{Spec}}}
\def\Hom{\mathop{\mathrm{Hom}}\nolimits}
\def\Ext{\mathop{\mathrm{Ext}}\nolimits}
\def\SheafHom{\mathop{\mathcal{H}\!\mathit{om}}\nolimits}
\def\SheafExt{\mathop{\mathcal{E}\!\mathit{xt}}\nolimits}
\def\Sch{\mathit{Sch}}
\def\Mor{\mathop{\mathrm{Mor}}\nolimits}
\def\Ob{\mathop{\mathrm{Ob}}\nolimits}
\def\Sh{\mathop{\mathit{Sh}}\nolimits}
\def\NL{\mathop{N\!L}\nolimits}
\def\CH{\mathop{\mathrm{CH}}\nolimits}
\def\proetale{{pro\text{-}\acute{e}tale}}
\def\etale{{\acute{e}tale}}
\def\QCoh{\mathit{QCoh}}
\def\Ker{\mathop{\mathrm{Ker}}}
\def\Im{\mathop{\mathrm{Im}}}
\def\Coker{\mathop{\mathrm{Coker}}}
\def\Coim{\mathop{\mathrm{Coim}}}

% Boxtimes
%
\DeclareMathSymbol{\boxtimes}{\mathbin}{AMSa}{"02}

%
% Macros for moduli stacks/spaces
%
\def\QCohstack{\mathcal{QC}\!\mathit{oh}}
\def\Cohstack{\mathcal{C}\!\mathit{oh}}
\def\Spacesstack{\mathcal{S}\!\mathit{paces}}
\def\Quotfunctor{\mathrm{Quot}}
\def\Hilbfunctor{\mathrm{Hilb}}
\def\Curvesstack{\mathcal{C}\!\mathit{urves}}
\def\Polarizedstack{\mathcal{P}\!\mathit{olarized}}
\def\Complexesstack{\mathcal{C}\!\mathit{omplexes}}
% \Pic is the operator that assigns to X its picard group, usage \Pic(X)
% \Picardstack_{X/B} denotes the Picard stack of X over B
% \Picardfunctor_{X/B} denotes the Picard functor of X over B
\def\Pic{\mathop{\mathrm{Pic}}\nolimits}
\def\Picardstack{\mathcal{P}\!\mathit{ic}}
\def\Picardfunctor{\mathrm{Pic}}
\def\Deformationcategory{\mathcal{D}\!\mathit{ef}}
\usepackage{cleveref}
\crefname{originallemma}{original lemma}{original lemmas}
\Crefname{originallemma}{Original lemma}{Original lemmas}
\makeatletter
\newlabel{derived-lemma-unbounded-right-derived}{{07K7}{source}{}{sourcepointer-1}{}}
\newlabel{derived-lemma-unbounded-right-derived@cref}{{[originallemma][0][]07K7}{[1][0][]source}{}{}{}}
\makeatother
\crefname{originallemma}{original lemma}{original lemmas}
\Crefname{originallemma}{Original lemma}{Original lemmas}
\makeatletter
\newlabel{injectives-lemma-derived-products}{{07D9}{source}{}{sourcepointer-2}{}}
\newlabel{injectives-lemma-derived-products@cref}{{[originallemma][0][]07D9}{[1][0][]source}{}{}{}}
\makeatother
\crefname{originaltheorem}{original theorem}{original theorems}
\Crefname{originaltheorem}{Original theorem}{Original theorems}
\makeatletter
\newlabel{injectives-theorem-K-injective-embedding-grothendieck}{{079P}{source}{}{sourcepointer-3}{}}
\newlabel{injectives-theorem-K-injective-embedding-grothendieck@cref}{{[originaltheorem][0][]079P}{[1][0][]source}{}{}{}}
\makeatother
\title{Bounded cohomological dimension and commutation with direct sums make the extension by zero of the structure sheaf compact}
\author{The Stacks Project Authors}
\date{stacks-project:sites-cohomology.tex:local-snapshot-20261011; GNU FDL 1.2 or later}
\begin{document}
\maketitle
\noindent\textit{Editorial note.} All delivered mathematical text is the author's own, copied from the pinned Stacks Project chapter file sites-cohomology.tex; no statement or proof step was rewritten, reordered or removed.\par\medskip
\noindent\textit{Editorial note.} The counted proof is the author's complete argument: the proof is written out in full: the reduction to the derived global sections, the use of the two hypotheses in the form of finite cohomological dimension and acyclicity of direct sums of injectives, the choice of K-injective representatives, and the computation on the resulting complex of acyclics.\par\medskip
\noindent\textit{Editorial note.} The label lemma-when-jshriek-compact names the counted statement, and the proof environment immediately following it in the source is delivered verbatim as its proof.\par\medskip
\noindent\textit{Editorial note.} The counted statement is carried by the author's own lemma environment in the chapter source.\par\medskip
\noindent\textit{Editorial note.} The statement and proof are copied from the pinned chapter file; the three cited inputs are kept as external pointers with their canonical Stacks tags.\par\medskip
\noindent\textit{Editorial note.} The preamble is the chapter's own retained minimal preamble (sources/stacks-sites-cohomology/retained-preamble.tex), a verbatim excerpt of the native preamble: the Xy-pic setup and the native theorem declarations with the mathematical macros.\par\medskip
\noindent\textit{Allowed standard prerequisites (editorial).}\begin{itemize}
\item Compactness of j\_!O\_U and the identification of its Hom functor. Same chapter, the identification Hom(j\_!O\_U, K) = R Gamma(U, K) and the notion of a compact object of the derived category. Supplies the reduction of the compactness statement to the commutation of R Gamma(U, -) with direct sums.
\item K-injective resolutions and the computation of right derived functors on acyclics. Injectives, the theorem providing a K-injective embedding, the lemma that direct sums of K-injectives represent direct sums, and Derived Categories, the lemma computing an unbounded right derived functor on a complex of acyclics. Supplies the representatives used to compute R Gamma(U, direct sum of K\_i).
\end{itemize}\medskip
\ifdefined\hypertarget\hypertarget{sourcepointer-1}{}\fi
\noindent\textit{Source reference.} Original lemma 07K7: sources/stacks-chow/tags.txt, source lines 9647--9647. Prerequisite pointer invoked at proof line(s) 14158; it names a standard fact of a sibling chapter that the counted argument uses, and every step of the counted proof built on it is retained verbatim in the delivered body.\par\medskip
\ifdefined\hypertarget\hypertarget{sourcepointer-2}{}\fi
\noindent\textit{Source reference.} Original lemma 07D9: sources/stacks-chow/tags.txt, source lines 9404--9404. Prerequisite pointer invoked at proof line(s) 14160; it names a standard fact of a sibling chapter that the counted argument uses, and every step of the counted proof built on it is retained verbatim in the delivered body.\par\medskip
\ifdefined\hypertarget\hypertarget{sourcepointer-3}{}\fi
\noindent\textit{Source reference.} Original theorem 079P: sources/stacks-chow/tags.txt, source lines 9279--9279. Prerequisite pointer invoked at proof line(s) 14156; it names a standard fact of a sibling chapter that the counted argument uses, and every step of the counted proof built on it is retained verbatim in the delivered body.\par\medskip
\begin{lemma}
\label{lemma-when-jshriek-compact}
Let $(\mathcal{C}, \mathcal{O})$ be a ringed site. Let $U$ be an object of
$\mathcal{C}$. The $\mathcal{O}$-module $j_!\mathcal{O}_U$ is a
compact object of $D(\mathcal{O})$ if there exists an integer $d$ such that
\begin{enumerate}
\item $H^p(U, \mathcal{F}) = 0$ for all $p > d$, and
\item the functors $\mathcal{F} \mapsto H^p(U, \mathcal{F})$
commute with direct sums.
\end{enumerate}
\end{lemma}

\begin{proof}
Assume (1) and (2). Recall that $\Hom(j_!\mathcal{O}_U, K) = R\Gamma(U, K)$ for
$K$ in $D(\mathcal{O})$. Thus we have to show that $R\Gamma(U, -)$
commutes with direct sums. The first assumption means that the functor
$F = H^0(U, -)$ has finite cohomological dimension. Moreover, the second
assumption implies any direct sum of injective modules is acyclic for $F$.
Let $K_i$ be a family of objects of $D(\mathcal{O})$.
Choose K-injective representatives $I_i^\bullet$ with injective terms
representing $K_i$, see Injectives, Theorem
\ref{injectives-theorem-K-injective-embedding-grothendieck}.
Since we may compute $RF$ by applying $F$ to any complex of acyclics
(Derived Categories, Lemma \ref{derived-lemma-unbounded-right-derived})
and since $\bigoplus K_i$ is represented by $\bigoplus I_i^\bullet$
(Injectives, Lemma \ref{injectives-lemma-derived-products})
we conclude that $R\Gamma(U, \bigoplus K_i)$ is represented by
$\bigoplus H^0(U, I_i^\bullet)$. Hence $R\Gamma(U, -)$ commutes
with direct sums as desired.
\end{proof}

\end{document}
